\section{Rare word fields, exact marks and the Poisson staircase}
\label{sec:patterns-fields}
The two-window estimate controls dependence between distant words. Their
local dependence has a different origin: a suffix of one word may coincide
with a prefix of another. Keeping these two issues separate gives a joint
law for a growing dictionary of prescribed words, and then a sign-resolved
law for long runs.

\subsection{A dictionary of prescribed words}\label{sec:word-dictionaries}
Joint word and clump approximations are classical in Markov sequences
\cite{ReinertSchbath1998}. Here the local overlap condition is kept
separate from the arithmetic estimate for separated windows.
Let $\mathcal W\subset\{\pm1\}^B$ be a nonempty deterministic set of distinct
words, and write $m=|\mathcal W|$, $p_B=2^{-B}$ and $\Lambda=Nmp_B$.
All signs, including the first one, are prescribed. Define
\[
 I_{x,w}=\1_{\{f(x-1+j)=w_j\ (0\le j<B)\}},
              \qquad x\in I_N,\quad w\in\mathcal W.
\]
The first letter is at $x-1$; this indexing matches the left-border
convention for runs.
For $1\le d<B$ write $w\mathrel{\sim_d}v$ when
$w_{d+j}=v_j$ for $0\le j<B-d$. The directed overlap weight is
\begin{equation}\label{eq:dictionary-overlap}
 \Omega(\mathcal W)=\frac1m\sum_{w,v\in\mathcal W}
                   \sum_{d=1}^{B-1}2^{-d}\1_{\{w\sim_d v\}}.
\end{equation}
The normalization makes $\Lambda\Omega$ the expected scale of ordered
local pairs, up to an absolute factor. It includes self-overlaps and
cross-overlaps; controlling each word separately would miss the latter.
For example, $(+,+,-)$ followed one site later by $(+,-,+)$ is compatible:
the two occurrences fit into $(+,+,-,+)$. The overlap is a constraint on
the words themselves, before the values of the multiplicative function
are considered.

\Needspace{25\baselineskip}
\begin{theorem}[Joint Poisson approximation for a word dictionary]
\label{thm:word-field}
Fix $0<\beta_-<\beta_+<\infty$ and $0<\varepsilon<1/3$. Uniformly for
$\beta_-\log N\le B\le\beta_+\log N$ and all dictionaries $\mathcal W$,
\begin{align}
 &\EE\,\dTV\left(\Law((I_{x,w})_{x\in I_N,w\in\mathcal W}
                       \mid\mathcal F_{Y_N^\sharp}),
            \bigotimes_{x\in I_N,w\in\mathcal W}\Pois(p_B)\right)
       \notag\\
 &\quad\ll_{\beta_-,\beta_+,\varepsilon}
 \Lambda\Omega(\mathcal W)
 +\Lambda(1+\Lambda)e^{-V_N+o(\nu_N)}\notag\\
 &\qquad\quad
 +N^{-1/3+\varepsilon}
       \left(\Lambda+\Lambda^{11/6}m^{1/6}
                       +\Lambda^{4/3}m^{2/3}\right).
       \label{eq:word-field-rate}
\end{align}
The implied constant and the little-oh term are independent of the words
and of $m$. In particular, fix $0<\eta<1/2$ and $C>0$. If
\begin{equation}\label{eq:dictionary-critical}
 m_N\le N^{1/2-\eta},\qquad
 |B_N-\log_2(Nm_N)|\le C,\qquad
 \Omega(\mathcal W_N)\longrightarrow0,
\end{equation}
then the entire site-and-word-labelled field converges in total variation
to the independent Poisson field. Its error is
\begin{equation}\label{eq:dictionary-critical-rate}
 O_{C,\eta}\!\left(\Omega(\mathcal W_N)
                  +e^{-V_N+o(\nu_N)}+N^{-\eta/2}\right).
\end{equation}
Every deterministic restriction of the site--word index set, and every
measurable function of the field, satisfies the same bound.
\end{theorem}
\begin{proof}
Use $Y=Y_N^\sharp$ and remove starts whose $B$-vertex support contains a
$Y$-defective integer. Let $D_Y^{\rm val}$ denote this set and put
$\mathcal M_B^{\rm val}=\sum_{x\in I_N}(2^{m_B(x)}-1)$ and $a=mp_B$.
Distinct words at a site are mutually exclusive, so $0<a\le1$ and
$\Lambda=Na$. By \cref{cor:absolute-marginal}, deleting the actual and
target coordinates costs at most
\[
 a\mathcal M_B^{\rm val}+2a|D_Y^{\rm val}|.
\]
All retained word events have conditional marginal $p_B$. Join all words
at the same site and all words at sites with intersecting prime supports
above $Y$, adding edges for overlapping supports if necessary. After
conditioning this is an exact dependency graph. Its first Poisson term is
$O(a^2(NB+E_Y))$, where $E_Y$ counts ordered pairs of distinct sites
sharing a prime above $Y$.

For $y=x+d$ with $1\le d<B$, two specified words are incompatible unless
$w\sim_d v$. In the compatible case they prescribe all $B+d$ values in
the union. Every vertex of a good union has a private coordinate above
$Y$, because its diameter is less than $2B<Y$. Hence
\begin{equation}\label{eq:absolute-overlap-law}
 \PP(I_{x,w}I_{x+d,v}=1\mid\mathcal F_Y)
          =2^{-(B+d)}\1_{\{w\sim_d v\}}.
\end{equation}
Summing ordered local pairs costs at most $2\Lambda\Omega(\mathcal W)$.
Different words at the same site contribute zero to the second term.

For separated good sites put $I_x^{\mathcal W}=\sum_{w\in\mathcal W}I_{x,w}$.
This is an indicator, not a count with possible multiplicities, and
$\PP(I_x^{\mathcal W}=1)=a$. Fourier inversion on the valuation rows also gives
$\PP(I_x^{\mathcal W}I_y^{\mathcal W}=1)\le a^2 2^{\rho_B^{\rm val}(x,y)}$. Combining the two
bounds \emph{before} summing over sites yields
\begin{align}
 \sum_{w,v\in\mathcal W}\PP(I_{x,w}I_{y,v}=1)
 &=\PP(I_x^{\mathcal W}I_y^{\mathcal W}=1)\notag\\
 &\le \min\{a,a^2 2^{\rho_B^{\rm val}(x,y)}\}\notag\\
 &\le a^2\left(1+
       \min\{a^{-1},2^{\rho_B^{\rm val}(x,y)}-1\}\right).
       \label{eq:dictionary-saturation}
\end{align}
Only the error summands have been grouped: the Poisson theorem is still
applied to the full family indexed by $(x,w)$. No word or position
label is discarded. Its averaged second term over separated graph
edges is therefore at most
$a^2(E_Y+R_2^{\mathrm{val},[1/a]}(N,B))$. We have proved the finite
transfer inequality
\begin{align}
 \EE\,\dTV
 &\ll a(\mathcal M_B^{\rm val}+|D_Y^{\rm val}|)
       +a^2\left(NB+E_Y+R_2^{\mathrm{val},[1/a]}(N,B)\right)
       +\Lambda\Omega(\mathcal W).
       \label{eq:dictionary-finite-transfer}
\end{align}

Here $\mathcal M_B^{\rm val}\le N^{1/2+o(1)}$,
$|D_Y^{\rm val}|/N\le e^{-V_N+o(\nu_N)}$, and
$E_Y/N^2\le e^{-V_N+o(\nu_N)}+N^{-1+o(1)}$.
Use \cref{prop:capped-profile} with $T=N/\Lambda$ and
$Q_B=2^B=Nm/\Lambda\ge T$. Its three contributions, after multiplication
by $a^2$, are bounded by
\[
 N^{-1/3+\varepsilon}
 \left(\Lambda^{11/6}m^{1/6}
                +\Lambda^{4/3}m^{2/3}+\Lambda\right).
\]
The smaller polynomial errors are absorbed with a smaller intermediate
$\varepsilon$: for example, $m\ge\Lambda/N$ implies
$\Lambda^2/N\le
N^{-1/3}\Lambda^{4/3}m^{2/3}$. This proves
\eqref{eq:word-field-rate}.
Under \eqref{eq:dictionary-critical}, $\Lambda$ is bounded above and
below and $B/\log N$ stays in one fixed compact interval. The polynomial
term is $O(N^{-\eta/2})$ on taking $\varepsilon=\eta/6$.
Restriction and measurable contraction cannot increase total variation.
\end{proof}

The exponent $1/2$ is a sufficient bound, not an optimality claim. The
remaining restriction comes from the nonterminal term $NQ_B^{2/3}$:
after normalization it gives $N^{-1/3+\varepsilon}m^{2/3}$ at bounded
intensity. The terminal term has instead been capped by the marginal
probability. This distinction matters: a large raw relation count need
not produce a proportionally large probability of seeing a dictionary
word.

\begin{corollary}[Replacement by independent signs]\label{cor:iid-word-replacement}
Let $(\xi_n)$ be independent fair signs and let $I^{\rm iid}_{x,w}$ be the
same word indicators formed from them. Under
\eqref{eq:dictionary-critical},
\[
 \dTV\left(\Law((I_{x,w})_{x,w}),
                     \Law((I^{\rm iid}_{x,w})_{x,w})\right)\longrightarrow0.
\]
In particular the vector of counts of the individual words is jointly
close to independent Poisson variables of mean $N2^{-B}$, even when the
number of words tends to infinity.
\end{corollary}
\begin{proof}
For the independent sequence the graph needs only overlapping supports.
Equation \eqref{eq:absolute-overlap-law} is exact without conditioning,
and its process Poisson error is
$O(\Lambda\Omega(\mathcal W)+\Lambda^2B/N)$.
Compare both fields with the same independent Poisson field and apply
\cref{thm:word-field}. Summing sites is a measurable contraction.
\end{proof}

\begin{corollary}[Most dictionaries have small overlap]\label{cor:generic-dictionaries}
Let $B\ge1$ and $1\le m\le2^B$. Choose an $m$-element subset $\mathcal W$
uniformly among the subsets of $\{\pm1\}^B$.
This auxiliary choice is independent of $f$. Then
\begin{equation}\label{eq:generic-overlap}
 \EE_{\mathcal W}\Omega(\mathcal W)=\frac{m(B-1)}{2^B}.
\end{equation}
Under the first two conditions in \eqref{eq:dictionary-critical}, the
fraction of dictionaries with $\Omega(\mathcal W)>N^{-1/2}$ is
$O_{C,\eta}(\log N/\sqrt N)$. Every dictionary outside this exceptional
set has the field approximation \eqref{eq:dictionary-critical-rate}.
\end{corollary}
\begin{proof}
Put $Q=2^B$. At each proper displacement $1\le d<B$, there are
$2^d$ self-compatible words and $2^{B+d}$ ordered compatible pairs.
The number of distinct compatible pairs is therefore $2^d(Q-1)$.
A word belongs to the dictionary with probability $m/Q$; a distinct
ordered pair belongs with probability $m(m-1)/(Q(Q-1))$.
After multiplication by $2^{-d}/m$, the expected contribution of this
displacement is exactly
\[
 \frac{2^{-d}}m\left(2^d\frac mQ
       +2^d(Q-1)\frac{m(m-1)}{Q(Q-1)}\right)=\frac mQ.
\]
Sum over the $B-1$ displacements to get \eqref{eq:generic-overlap}.
This also covers $m=1$; for $B=1$ the sum is empty.
At criticality $m2^{-B}=\Lambda/N=O_C(1/N)$ and $B=O_\eta(\log N)$.
Markov's inequality proves the exceptional-fraction estimate. The result
is a statement about the size of a class of deterministic dictionaries;
no additional randomness is required in \cref{thm:word-field}.
\end{proof}

\begin{corollary}[Explicit dictionaries with no overlaps]
\label{cor:marker-dictionaries}
For all sufficiently large $B$, there is an explicitly described family
$\mathcal C_B\subset\{0,1\}^B$ of size at least $2^B/(32B)$ with no
proper suffix of any word equal to a prefix of any word, including itself.
Every subset of $\mathcal C_B$, after identifying $0,1$ with $-1,+1$,
has $\Omega=0$. In particular deterministic examples exist throughout
the range \eqref{eq:dictionary-critical}.
\end{corollary}
\begin{proof}
This is the binary marker construction of Chee, Kiah, Purkayastha and
Wang~\cite[Section III]{CheeKiahPurkayasthaWang2013}; we include the
short argument and the coarse size bound needed here.
Set $k=\lceil\log_2(2B)\rceil$ and take all words
\[
                0^k\,1\,u\,1,
       \qquad |u|=B-k-2,
       \qquad u\text{ contains no }0^k.
\]
The only occurrence of $0^k$ in such a word starts at its first position:
the first following letter and the last letter are both $1$.
An overlap of length at most $k$ is impossible because the prefix is all
zeros and the suffix ends in $1$. A longer proper overlap would put
$0^k$ at a later position, also impossible.
There are $2^{B-k-2}$ choices before the restriction on $u$. A union
bound over at most $B$ positions excludes at most the fraction
$B2^{-k}\le1/2$. Thus
$|\mathcal C_B|\ge2^{B-k-3}\ge2^B/(32B)$.
At the critical dictionary length, $m/2^B=O(1/N)$ and $B=O(\log N)$,
so any required $m$ words can be selected from this family for large $N$.
\end{proof}
For example, one may select $\lfloor N^{2/5}\rfloor$ such words at length
$B=(7/5)\log_2N+O(1)$. Their full labelled field has a Poisson
approximation. The coding construction is classical; the conclusion for
the multiplicative sequence is the application of the theorem above.

\subsection{Words up to a common sign}
For $a=(a_0,\ldots,a_L)\in\{\pm1\}^{L+1}$ put
$I_x(a)=I_{x,a}+I_{x,-a}$. Define
\begin{equation}\label{eq:overlap-definition}
 \mathcal O(a)=\{1\le d\le L:\exists\eta\in\{\pm1\},\
         a_{d+j}=\eta a_j\ (0\le j\le L-d)\},\qquad
 \Theta(a)=\sum_{d\in\mathcal O(a)}2^{-d}.
\end{equation}
For the two-word dictionary $\{a,-a\}$ one has
$\Omega(\{a,-a\})=\Theta(a)$: every compatible shift has two directed
sign choices.
\begin{corollary}[Poisson law for low-overlap patterns]\label{thm:patterns}
Fix $C>0$ and $\varepsilon>0$. Uniformly in $a=a_N$ and
$|L-\log_2N|\le C$,
\begin{equation}\label{eq:pattern-rate}
 \dTV\left(\Law\left(\sum_{x\in I_N}I_x(a)\right),\Pois(N2^{-L})\right)
 \ll_{C,\varepsilon}\Theta(a)+e^{-V_N+o_C(\nu_N)}+N^{-1/3+\varepsilon}.
\end{equation}
The same bound holds under any deterministic spatial mask, with target
mean $|A_N|2^{-L}$.
\end{corollary}
\begin{proof}
Apply \cref{thm:word-field} with $B=L+1$, $m=2$, and merge the two word
coordinates. Their events are disjoint and their target means add.
\end{proof}
If the least compatible shift $d_*(a_N)$ tends to infinity, then
$\Theta(a_N)\le2^{1-d_*(a_N)}\to0$. For the run-start word
$(-1,+1,\ldots,+1)$ only $d=L$ is compatible. Periodic words can instead
retain a short overlap and need declustering. Notice that the preceding
dictionary theorem also separates $a$ and $-a$: their asymptotic
independence is not assumed in this corollary.

The dictionary comparison now controls the rare occurrences of each
prescribed word. We return to constant runs and keep their attributes
together: a start, its full length, its sign and all thresholds it exceeds.

\subsection{Exact marks without a growing-mark loss}
Almost surely every rightward constant run is finite. Indeed constancy on
$[x,\infty)$ would make all prime signs beyond $x$ equal, an event of
probability zero; take a countable union over $x$. Thus
\begin{equation}\label{eq:mark-partition}
 J_{x,L}=\sum_{e\ge0}K_{x,e},\qquad
 K_{x,e}=\1_{\{f(x-1)=-f(x),\ f(x)=\cdots=f(x+L+e-1),\
                              f(x+L+e)=-f(x)\}}.
\end{equation}
This system has $L+e+1$ independent homogeneous rows on a good support and
conditional marginal $p_e=2^{-L-e-1}$. Its homogeneous rows are nested in
those for a larger mark, although the terminal affine right-hand side
changes. Marks at the same start are mutually exclusive. We may retain the sign as
well:
\begin{equation}\label{eq:signed-mark}
 K^s_{x,e}=K_{x,e}\1_{\{f(x)=s\}},\qquad
 p_{e,s}=2^{-L-e-2},\qquad s\in\{\pm1\}.
\end{equation}
On a good support this event prescribes every value, so its conditional
probability is $p_{e,s}$. The exact identities are
$\sum_sK^s_{x,e}=K_{x,e}$ and $\sum_{e,s}K^s_{x,e}=J_{x,L}$.

\begin{theorem}[Finite exact-mark comparison]\label{thm:marked-field}
Let $E\ge0$, $Q=L+E+1$, and assume
$\beta_-\log N\le L+1\le Q+1\le\beta_+\log N$.
For $Y>2(Q+1)$ let $D_Y(Q)$ be the starts whose support
$\{x-1,\ldots,x+Q-1\}$ contains a $Y$-defective integer, and let $E_Y(Q)$
be the ordered graph-edge count for these maximal supports. Then
\begin{align}
 &\EE\,\dTV\left(\Law((K_{x,e})_{x\in I_N,\ 0\le e\le E}\mid\mathcal F_Y),
                 \bigotimes_{x\in I_N,\ 0\le e\le E}\Pois(p_e)\right)
 \notag\\
 &\quad\ll p\mathcal M_B+p|D_Y(Q)|
       +p^2\bigl(NQ+E_Y(Q)+R_2(N,L)\bigr),\qquad p=2^{-L}.
       \label{eq:marked-finite}
\end{align}
Exactly the same upper bound holds for the sign-resolved family
$(K^s_{x,e})_{x\in I_N,\,0\le e\le E,\,s\in\{\pm1\}}$, compared with
independent Poisson variables of means $p_{e,s}$. There is no polynomial
factor in the number of marks. At $Y=Y_N^\sharp$,
the right side is at most
\begin{equation}\label{eq:marked-hard}
 C_\varepsilon\lambda_N(1+\lambda_N)
       \left(e^{-V_N+o(\nu_N)}+N^{-1/3+\varepsilon}\right).
\end{equation}
\end{theorem}
\begin{proof}
On good maximal supports all marginals are $p_e$. Apply the process
Poisson theorem to the marked family, allowing all marks on neighbouring
supports to be neighbours, and all marks at the same site to be neighbours.
The two identities
\begin{equation}\label{eq:mark-folding}
 \sum_{e=0}^Ep_e\le p,\qquad
 \sum_{e,f=0}^E\PP(K_{x,e}K_{y,f}=1)\le\PP(J_{x,L}J_{y,L}=1)
\end{equation}
remove any double sum over marks before estimating the pair mass.
Local pairs cost $O(NQp^2)$ and separated pairs cost
$p^2(E_Y(Q)+R_2(N,L))$. Bad actual marks at one site have total mass at
most that of $J_{x,L}$, so their deletion is bounded by
$p\mathcal M_B+p|D_Y(Q)|$, regardless of which late vertex caused removal.
This proves \eqref{eq:marked-finite}.
The bounds $|D_Y(Q)|\le(Q+1)A(3N,Y)$ and \eqref{eq:edge-count} with
$B$ replaced by $Q+1$, together with \cref{thm:master-profile,prop:cutoffs},
give \eqref{eq:marked-hard}.

For the signed version, the same-site alternatives are still disjoint and
both sums in \eqref{eq:mark-folding} remain valid after summing signs.
It remains to check local pairs before that summation. Put $a=L+e$,
$b=L+f$, and $d=y-x>0$. If $d<a$ the pair is impossible. If $d=a$, the
two prescribed value strings share two vertices, require the second run
sign to be the opposite of the first, and have probability
$4p_{e,s}p_{f,t}$ when compatible. If $d=a+1$, they share one vertex,
require equal run signs, and have probability $2p_{e,s}p_{f,t}$.
For $d>a+1$ their actual supports are disjoint, and a good local union has
probability $p_{e,s}p_{f,t}$. These assertions follow by counting the
prescribed values in a union whose private coordinates make them uniform.
Thus the local sum is still $O(NQp^2)$ and all other terms are unchanged.
\end{proof}
The use of \eqref{eq:mark-folding} is essential. An exact mark $e<E$ does
not imply a threshold event at length $L+E+1$; a bad vertex can occur beyond
its support. We bound deletion by the base event, not the longest event.

The exact tail identity is
\begin{equation}\label{eq:mark-tail}
 \{\exists x\in I_N,e>E:K_{x,e}=1\}=
                      \{Z_{N,L+E+1}>0\}.
\end{equation}
The first-moment estimate bounds this probability by
$\lambda_N2^{-(E+1)}(1+N^{-1/2+o(1)})$ whenever the enlarged length is in
a fixed logarithmic band. The independent target tail has mass
$\lambda_N2^{-(E+1)}$. Couple both tails to zero.
At criticality choose $E=\lceil V_N/\log2\rceil$. This is $o(\log N)$;
\eqref{eq:marked-hard} and \eqref{eq:mark-tail} prove the full-field
estimate, including signs, in \cref{thm:poisson-main}. Riemann sums and
geometric tail tightness give its weak diffuse limit.

\begin{corollary}[An explicit compound-Poisson cluster law]
\label{cor:window-clusters}
Let $C_{N,L}$ count all constant length-$L$ windows whose left endpoint
lies in $I_N$, without requiring left maximality. At criticality its law
is at distance $o(1)$ from that of
\[
 \sum_{i=1}^{P_N}H_i,\qquad P_N\sim\Pois(\lambda_N),\qquad
 \PP(H_i=h)=2^{-h},\quad h\ge1,
\]
with independent $H_i$ and $P_N$. The target probability generating function
is $\exp[\lambda_N(z/(2-z)-1)]$ for $0\le z\le1$.
\end{corollary}
\begin{proof}
A maximal run of length $L+e$ contributes $e+1$ windows. Thus, apart from
runs crossing an endpoint of $I_N$,
$C_{N,L}=\sum_{x\in I_N,e\ge0}(e+1)K_{x,e}$.
A discrepancy implies a constant length-$L$ window starting at $N$ or
at $2N-1$. Each such event has $L-1$ relative constraints; private pivots
and the pointwise defect bound give probability at most
$2^{-(L-1)+O(\log N/\log\log N)}=N^{-1+o(1)}$.
The full marked-field comparison contracts under the measurable weighted
sum. Its target is the stated compound Poisson law, since $H=e+1$.
No moment convergence is inferred from this contraction.
\end{proof}

\subsection{Growing depth: distinguish positions from counts}
Put $b_N=\lfloor\log_2N\rfloor$, $\vartheta_N=\log_2N-b_N$. For
$d_N\in\NN$ with $d_N=o(\log N)$, set
\[
 L_N=b_N-d_N,\qquad \Lambda_N=N2^{-L_N}=2^{\vartheta_N+d_N}.
\]
Let $K^{\rm ex}_{x,q}$ denote the exact-length event for a left-maximal run
of length $q$ (zero for $q<1$). The moving field and its target are
\begin{equation}\label{eq:moving-field}
 \mathcal E_{N,d_N}=\sum_{x\in I_N}\sum_{r\ge-d_N}
             K^{\rm ex}_{x,b_N+r}\delta_{(x/N,r)},\qquad
 \EE\Pi_{N,d_N}\{(x/N,r)\}=2^{-b_N-r-1},
\end{equation}
where the target atoms are independent Poisson variables. Their total mean
is $\Lambda_N$. Define exact-level and threshold counts by
\[
 X_{N,r}=\sum_{x\in I_N}K^{\rm ex}_{x,b_N+r},\qquad
 T_{N,m}=\sum_{r\ge m}X_{N,r}=Z_{N,b_N+m}.
\]
For the target write $P_{N,r}=\Pi_{N,d_N}([1,2]\times\{r\})$ and
$S_{N,m}=\sum_{r\ge m}P_{N,r}$. Thus
$P_{N,r}\sim\Pois(2^{\vartheta_N-r-1})$ independently.

The preceding estimates bound the mean, over the small-prime environment,
of a conditional total-variation distance. Restricting this average to an
event $C_N\in\mathcal F_Y$ multiplies the error by at most
$1/\PP(C_N)$; this is the stable-kernel principle of
\cref{lem:stable-lift}. It is the first information cost below.

\begin{theorem}[Two levels of joint approximation]\label{thm:moving-comparison}
Fix $c>0$ and $0<\varepsilon<1/3$. Let $C_N$ be an arithmetic
conditioning event of positive probability and
put $I_N^{\rm info}=-\log\PP(C_N)$.

\textup{(i) Labelled field.} If $C_N\in\mathcal F_{Y_N^\sharp}$, a
sufficient condition for full-field total variation to tend to zero is
\begin{equation}\label{eq:labelled-hard-budget}
 I_N^{\rm info}+\log\Lambda_N+\log(1+\Lambda_N)\le V_N-c\nu_N.
\end{equation}
At a movable cutoff, if $C_N\in\mathcal F_{\widehat Y_N}$, the sufficient
condition improves to
\begin{equation}\label{eq:labelled-budget}
 I_N^{\rm info}+\log\Lambda_N\le\widehat V_N/2-c\widehat\nu_N.
\end{equation}
In the nontrivial range, the latter comparison has upper bound
\begin{align}
 &e^{I_N^{\rm info}+\log\Lambda_N-\widehat V_N/2+o(\widehat\nu_N)}
 +e^{I_N^{\rm info}+2\log\Lambda_N-\widehat V_N+o(\widehat\nu_N)}\notag\\
 &\hspace{20mm}+e^{I_N^{\rm info}}\Lambda_N(1+\Lambda_N)
                         N^{-1/3+\varepsilon+o(1)}.
 \label{eq:labelled-rate}
\end{align}
\textup{(ii) Aggregated exact levels.} If $C_N\in\mathcal F_{Y_N^\sharp}$
and
\begin{equation}\label{eq:staircase-budget}
 I_N^{\rm info}+\log\Lambda_N\le V_N-c\nu_N,
\end{equation}
then, for every $0<c'<c$ and $0<\varepsilon<1/3$,
\begin{equation}\label{eq:staircase-rate}
 \dTV\left(\Law((X_{N,r})_{r\ge-d_N}\mid C_N),
       \bigotimes_{r\ge-d_N}\Pois(2^{\vartheta_N-r-1})\right)
 \ll e^{-c'\nu_N}+N^{-1/3+\varepsilon}.
\end{equation}
Exactly the same distance holds for the full threshold paths
$(T_{N,m})_{m\ge-d_N}$ and $(S_{N,m})_{m\ge-d_N}$.
\end{theorem}
\begin{proof}
Conditioning multiplies any averaged kernel error by at most
$e^{I_N^{\rm info}}$. For (i), apply \cref{thm:marked-field} with base
$L_N$ and then remove the mark tail using \eqref{eq:mark-tail}.
For the hard bound one may take
$E=\lceil(I_N^{\rm info}+V_N+\log(2+\Lambda_N))/\log2\rceil$.
Under \eqref{eq:labelled-hard-budget} this is $o(\log N)$ and both tail
costs are absorbed by the displayed error.
At a free cutoff $Y=e^w$, the leading process errors are
\[
 e^{I_N^{\rm info}}\Lambda_N e^{-D(H_N/w)+o(H_N/w)},\qquad
 e^{I_N^{\rm info}}\Lambda_N^2e^{-w+o(H_N/w)}.
\]
Choose $w=\widehat V_N$, so $D(H_N/w)=w/2$, and take an equally large
$o(\log N)$ tail cutoff. This proves \eqref{eq:labelled-rate} and
\eqref{eq:labelled-budget}. The hard choice remains sharper at bounded
intensity; the movable choice improves the admissible growing intensity.

For (ii), aggregate sites \emph{before} solving the multivariate Poisson
Stein equation. If $q_e=2^{-(e+1)}$ are the normalized category masses,
Barbour's Lemmas~2--3 \cite[p.~179]{Barbour1988SteinProcess}, as
applied in the companion, bound the directional second-difference factor by
\[
 C\frac{1+\log^+(2\Lambda_N)}{\Lambda_N\sqrt{q_eq_f}}.
\]
The mark sum remains uniform because
$\sum_{e\ge0}\sqrt{q_e}=1+\sqrt2$.
The full argument and the precise finite-cutoff bound are given in
\cref{supp:thm:aggregated}. It yields a leading error
$e^{I_N^{\rm info}}\Lambda_N(1+\log^+(2\Lambda_N))e^{-V_N+o(\nu_N)}$
and a polynomial remainder after the enlarged-length arithmetic estimate.
All subpolynomial factors in that remainder are absorbed by choosing a
smaller initial $\varepsilon$. Choose a constant strictly between $c'$ and
$c$ for the tail cutoff. The scalar relative estimate at the shifted length
bounds both tails by $O(e^{-c'\nu_N})$; see \cref{supp:thm:aggregated}.
Finally, on finite-mass sequences, the map $X_r\mapsto\sum_{s\ge r}X_s$
is a measurable bijection with inverse $X_r=T_r-T_{r+1}$. Total variation
is preserved, not merely contracted.
\end{proof}
Part (ii) retains every statistic of the multiset of lengths, but not the
positions. A maximal coupling identifies all those retained statistics at
once, with no union bound over thresholds. A field coupling from (i) also
retains positions. The ranges are sufficient bounds for these methods, not
phase transitions. In particular the full one-factor labelled range is not
proved by these comparisons.


\begin{corollary}[Sign-resolved growing fields]\label{cor:signed-growing}
Both conclusions of \cref{thm:moving-comparison} remain valid after retaining
the sign of each run. The labelled target has atom means $2^{-b_N-r-2}$
for each $s\in\{\pm1\}$, and the aggregated target has independent means
$2^{\vartheta_N-r-2}$ at each pair $(r,s)$.
\end{corollary}
\begin{proof}
For the labelled result use the signed version of \cref{thm:marked-field}
and the unchanged tail identity. For aggregation, the normalized category
weights are $q_{e,s}=2^{-e-2}$ and
$\sum_{e,s}\sqrt{q_{e,s}}=\sqrt2(1+\sqrt2)$.
The directional Stein bound therefore has the same summability.
At the largest retained support, the full value-relation space dominates
each mixed signed kernel by extension of coefficient vectors by zero;
\cref{prop:absolute-profile} supplies the same two-window envelope.
The retained support has $L+E+2$ vertices, differing by only one from the
relative-mark convention. Since $E=O(V_N)=o(\log N)$, its polynomial
remainder is absorbed exactly as in part (ii). The signed local-pair calculation and the same tail bound finish
the proof.
\end{proof}

\subsection{Target kernels, extremes and the Gaussian lower edge}
The independent target levels give, exactly,
\begin{align}
 S_{N,m}&\sim\Pois(2^{\vartheta_N-m}),\label{eq:target-marginal}\\
 S_{N,m+j}\mid S_{N,m}=n&\sim\operatorname{Bin}(n,2^{-j}),\label{eq:target-forward}\\
 S_{N,m-j}\mid S_{N,m}=n&\ \stackrel d=\ n+
      \Pois(2^{\vartheta_N-m}(2^j-1)),\label{eq:target-reverse}
\end{align}
when the lower levels are retained. The whole forward path is the number
of survivors among $n$ independent geometric lifetimes, each with
$\PP(E=e)=2^{-(e+1)}$. These are identities for the Poisson target, not
exact finite-$N$ Markov assertions for the multiplicative sequence.

For any integer $m\ge0$, the dyadic excess maximum
$M_{N,L}=\max\{\ell_x-L:x\in I_N,J_{x,L}=1\}$, with empty value $-\infty$,
satisfies the exact identity
\begin{equation}\label{eq:excess-void}
 \{M_{N,L}\le m\}=\{Z_{N,L+m+1}=0\}.
\end{equation}
Thus the scalar estimate at $L'=L+m+1$, with $\mu=N2^{-L'}$, gives
\[
 \PP(M_{N,L}\le m)=e^{-\mu}
 +O\left(\min\{1,\mu(e^{-V_N+o(\nu_N)}+N^{-1/3+\varepsilon})\}\right)
\]
whenever $L'+1$ is in a fixed logarithmic band. In particular one-hit upper
tails are relative when $N2^{-L'}\to0$.
For the $k$th largest run length among starts in $I_N$, along a subsequence
$\vartheta_N\to\vartheta$,
\begin{equation}\label{eq:order-statistic}
 \PP(R_N^{[k]}-b_N<m)\longrightarrow
 e^{-2^{\vartheta-m}}\sum_{j=0}^{k-1}\frac{(2^{\vartheta-m})^j}{j!}.
\end{equation}
This is just the event that fewer than $k$ starts exceed the threshold.

\begin{theorem}[Joint Poisson--Gaussian limit]\label{thm:bridge}
Assume $d_N\to\infty$ and the hypotheses of part (ii) of
\cref{thm:moving-comparison}. Put
\[
 G_N(j)=\frac{Z_{N,b_N-d_N+j}-\Lambda_N2^{-j}}{\sqrt{\Lambda_N}}.
\]
Along $\vartheta_N\to\vartheta$, for every fixed $J$ and every finite
$R\subset\Z$, conditionally on $C_N$,
\[
 ((G_N(j))_{0\le j\le J},(X_{N,r})_{r\in R})
 \Longrightarrow ((G(j))_{0\le j\le J},(P_r^{(\vartheta)})_{r\in R}),
\]
where $\EE G(j)G(k)=2^{-\max(j,k)}$, the variables
$P_r^{(\vartheta)}\sim\Pois(2^{\vartheta-r-1})$ are independent, and the
Gaussian and Poisson vectors are independent. The normalized Gaussian
sequence $H(j)=2^{j/2}G(j)$ is the stationary AR(1) chain
$H(j+1)=2^{-1/2}H(j)+2^{-1/2}\zeta_{j+1}$,
where $(\zeta_j)_{j\ge1}$ are independent standard normal random variables,
independent of $H(0)\sim N(0,1)$.
\end{theorem}
\begin{proof}
Transfer to the independent target by \eqref{eq:staircase-rate}.
Choose a fixed $r_0<\min R$ (any fixed $r_0$ if $R$ is empty).
The lower-edge sums restricted to $r<r_0$ are independent of the displayed
critical levels and satisfy the multivariate Poisson central limit theorem.
Their covariance after normalization tends to $2^{-\max(j,k)}$.
The centered sum over $r\ge r_0$ has bounded variance and hence vanishes
after division by $\sqrt{\Lambda_N}$. This proves the joint limit and
independence. The covariance of $H$ is $2^{-|j-k|/2}$, giving the stated
Gaussian autoregression.
\end{proof}
For fixed compact sets of levels, the untruncated centered field has the
phase-indexed target
\[
 {\rm PPP}\left(dt\otimes\sum_{r\in\Z}2^{\vartheta_N-r-1}\delta_r\right).
\]
Riemann sums give the corresponding Laplace approximation without assuming
that $\vartheta_N$ converges. A phase-convergent subsequence turns it into
a fixed weak limit; its intensity is locally finite although it diverges
towards $r=-\infty$. Full formulas and conditional configuration statements
are in \cref{supp:sec:target-consequences}. No continuous run lengths or
canonical coupling simultaneously over all $N$ is asserted.
